\documentclass[11pt,a4paper]{article}
\usepackage[utf8]{inputenc}
\usepackage[vietnamese]{babel}
\usepackage{amsmath,amssymb}
\usepackage{geometry}
\geometry{a4paper, total={180mm,260mm}, left=15mm, right=15mm, top=18mm, bottom=18mm}
\usepackage{tikz}
\usetikzlibrary{shapes.geometric, arrows.meta, positioning, calc}
\usepackage{tcolorbox}
\tcbuselibrary{skins,breakable}
\usepackage{xcolor}
\usepackage{multicol}
\usepackage{array}
\usepackage{fontawesome}
\usepackage{fancyhdr}

\pagestyle{fancy}
\fancyhf{}
\renewcommand{\headrulewidth}{0.5pt}
\renewcommand{\footrulewidth}{0.5pt}
\fancyhead[L]{\small\bfseries\color{teal}\faBook\ SỔ TAY KIẾN THỨC TOÁN 3}
\fancyhead[R]{\small\bfseries\color{teal}Lớp Toán Thầy Lĩnh -- loptoanthaylinh.vn}
\fancyfoot[C]{\small\color{gray}-- Sổ tay kiến thức: Bài 72. Luyện tập chung: Phép nhân và phép chia phạm vi 100 000 --}

\definecolor{primary}{HTML}{005F73}
\definecolor{secondary}{HTML}{0A9396}
\definecolor{accent}{HTML}{EE9B00}
\definecolor{danger}{HTML}{AE2012}
\definecolor{boxbg}{HTML}{F4F6F8}

\begin{document}

% --- BANNER TIÊU ĐỀ ---
\begin{tcolorbox}[
  enhanced,
  colback=white,
  colframe=primary,
  arc=3mm,
  boxrule=1.2pt,
  left=4mm, right=4mm, top=2.5mm, bottom=2.5mm,
  drop shadow=black!10!white
]
  \centering
  {\Large\bfseries\color{primary}\faBookmark\ BÀI 72: LUYỆN TẬP CHUNG: NHÂN CHIA PHẠM VI 100 000}\\[1.5mm]
  {\small\itshape\color{black!75}Đại tổng kết đặt tính nhân chia số có 5 chữ số, tính giá trị biểu thức và giải bài toán rút về đơn vị trên 1 trang A4}
\end{tcolorbox}

\vspace{-2mm}


\noindent
\begin{minipage}[t]{0.49\textwidth}\vspace{0pt}
\begin{tcolorbox}[
  enhanced,
  colback=blue!4!white,
  colframe=blue!75!black,
  arc=2mm, boxrule=0.8pt,
  left=2mm, right=2mm, top=1.8mm, bottom=1.8mm,
  title={\footnotesize\bfseries\color{white}\faCalculator\ 1. Kĩ thuật nhân chia 5 chữ số}
]
{\footnotesize
\begin{itemize}\setlength{\itemsep}{1.5pt}\setlength{\leftskip}{-2.5mm}
  \item Nhân số có 5 chữ số: nhân phải sang trái, cộng số nhớ vào hàng liền kề.
  \item Chia số có 5 chữ số: chia trái sang phải, không được quên chữ số $0$ ở thương khi số bị chia bé hơn số chia.
  \item Tính nhẩm tròn chục nghìn:
  $$30\,000 \times 3 = 90\,000 \qquad 80\,000 : 2 = 40\,000$$
\end{itemize}
}
\end{tcolorbox}
\end{minipage}
\hfill
\begin{minipage}[t]{0.49\textwidth}\vspace{0pt}
\begin{tcolorbox}[
  enhanced,
  colback=teal!4!white,
  colframe=teal!75!black,
  arc=2mm, boxrule=0.8pt,
  left=2mm, right=2mm, top=1.8mm, bottom=1.8mm,
  title={\footnotesize\bfseries\color{white}\faSliders\ 2. Bài toán giải bằng 2 bước tính}
]
{\footnotesize
\begin{itemize}\setlength{\itemsep}{1.5pt}\setlength{\leftskip}{-2.5mm}
  \item \textbf{Bài toán Rút về đơn vị:}
  \begin{itemize}\setlength{\itemsep}{0.5pt}\setlength{\leftskip}{-2.5mm}
    \item Bước 1: Tìm giá trị của $1$ đơn vị (Làm phép \textbf{CHIA}).
    \item Bước 2: Tìm giá trị của nhiều đơn vị đó (Làm phép \textbf{NHÂN}).
  \end{itemize}
  \item \textbf{Thứ tự thực hiện biểu thức số:}
  $$\text{Ngoặc trước} \to \text{Nhân / Chia} \to \text{Cộng / Trừ}$$
\end{itemize}
}
\end{tcolorbox}
\end{minipage}

\vspace{-2.5mm}
% --- KHỐI 3: SƠ ĐỒ TƯ DUY TIKZ ---
\noindent
\begin{tcolorbox}[
  enhanced,
  colback=orange!4!white,
  colframe=orange!80!black,
  arc=2mm, boxrule=0.8pt,
  left=2mm, right=2mm, top=1.5mm, bottom=1.5mm,
  title={\footnotesize\bfseries\color{white}\faSitemap\ SƠ ĐỒ TƯ DUY: HAI BƯỚC BÀI TOÁN RÚT VỀ ĐƠN VỊ}
]
\centering
\begin{tikzpicture}[
  box/.style={draw, rounded corners=2pt, align=center, font=\scriptsize\bfseries, inner sep=2.5pt, minimum height=0.75cm},
  arr/.style={-{Stealth[length=2mm]}, thick}
]
  \node[box, fill=purple!15, draw=purple!80!black, text width=3.2cm] (r1) at (0, 0) {ĐỀ BÀI CHO\\Giá trị của $n$ phần bằng nhau};
  \node[box, fill=teal!15, draw=teal!80!black, text width=3.4cm] (r2) at (4.6, 0) {BƯỚC 1: RÚT VỀ $1$ ĐƠN VỊ\\Lấy tổng chia cho $n$\\($A : n \implies$ Giá trị $1$ phần)};
  \node[box, fill=orange!15, draw=orange!80!black, text width=3.2cm] (r3) at (9.4, 0) {BƯỚC 2: TÍNH $k$ PHẦN\\Lấy giá trị $1$ phần $\times k$\\Đáp số bài toán};

  \draw[arr, purple!80!black] (r1) -- (r2);
  \draw[arr, teal!80!black] (r2) -- (r3);
\end{tikzpicture}
\end{tcolorbox}

\vspace{-2.5mm}
% --- KHỐI 4: VÍ DỤ & MẸO TRÁNH BẪY ---
\noindent
\begin{minipage}[t]{0.49\textwidth}\vspace{0pt}
\begin{tcolorbox}[
  enhanced,
  colback=green!4!white,
  colframe=green!60!black,
  arc=2mm, boxrule=0.8pt,
  left=2mm, right=2mm, top=1.8mm, bottom=1.8mm,
  title={\footnotesize\bfseries\color{white}\faCheckCircle\ Ví dụ thực hành chuẩn mực}
]
{\footnotesize
\textbf{Bài toán Rút về đơn vị:}\\
Mua $4$ quyển vở cùng loại hết $24\,000$ đồng. Hỏi mua $7$ quyển vở như thế hết bao nhiêu tiền?\\
\textbf{Bài giải chuẩn:}\\
Giá tiền $1$ quyển vở là:
$$24\,000 : 4 = 6000\text{ (đồng)}$$
Mua $7$ quyển vở như thế hết số tiền là:
$$6000 \times 7 = \mathbf{42\,000}\text{ (đồng)}$$
\phantom{Bài giải: }\textbf{Đáp số:} $42\,000$ đồng.
}
\end{tcolorbox}
\end{minipage}
\hfill
\begin{minipage}[t]{0.49\textwidth}\vspace{0pt}
\begin{tcolorbox}[
  enhanced,
  colback=red!4!white,
  colframe=red!70!black,
  arc=2mm, boxrule=0.8pt,
  left=2mm, right=2mm, top=1.8mm, bottom=1.8mm,
  title={\footnotesize\bfseries\color{white}\faExclamationTriangle\ Mẹo tránh bẫy trắc nghiệm}
]
{\footnotesize
\begin{itemize}\setlength{\itemsep}{1.5pt}\setlength{\leftskip}{-2.5mm}
  \item \textbf{Bẫy làm phép nhân trước phép chia trong bài toán rút về đơn vị:}\\
  Không thể lấy $24\,000 \times 7$ rồi mới chia $4$ vì phép tính số lớn rất dễ nhầm lẫn. Luôn tìm giá trị của \textbf{$1$ đơn vị} trước bằng phép chia!
  \item Ghi đúng tên đơn vị trong ngoặc (đồng, quyển, kg...).
\end{itemize}
}
\end{tcolorbox}
\end{minipage}


\end{document}
